\section{Introduction}\label{sec:intro}

Finite SAT encodings are useful for locating sharp feasibility boundaries in
social choice.  They can expose which combinations of axioms are inconsistent,
which assumptions form a minimal conflict, and which removals restore
satisfiability.  In an auditable possibility atlas, however, a solver result is
only one layer of the evidence.  The semantic statement, the encoding, the
solver status, the proof artifact, and the human-facing repair explanation have
different invariance requirements.

This paper studies one such requirement: whether a repair explanation remains
stable after a finite impossibility witness is transferred to a neighboring
case.  The motivating base case is the Sen liberal paradox with two voters and
four alternatives \citep{sen1970}.  In the M1.5 study, a bundled
\texttt{minlib} lever and a split presentation using
\texttt{decisive\_voter0} and \texttt{decisive\_voter1} generated the same
clause multiset under a verified variable renaming.  Nevertheless, their raw
lever-level inclusion-minimal repairs differed.  The bundled presentation
reported \lever{minlib}; the split presentation reported the two
person-specific singleton repairs separately.  Clause-multiset equivalence
ruled out a stronger-encoding explanation for that divergence.

M2.1 asks whether this repair-presentation witness transfers to larger finite
cases.  This is not the same question as whether the underlying impossibility
status transfers.  The repository's M2 layer owns the semantic transfer claim;
M2.1 operates one layer below it and asks what happens to the representation of
repair candidates.  The distinction matters because status can remain UNSAT
even when a split encoding introduces new selector structure, and because a
repair divergence is strongest when the compared clause multisets are identical
up to renaming.

We examine the two nearest directions away from $(2,4)$:
\begin{enumerate}
  \item increase the number of alternatives while fixing the two-voter witness
  structure, using $(2,5)$; and
  \item increase the number of voters while retaining four alternatives, using
  $(3,4)$.
\end{enumerate}
These directions require different split representations.  With two voters,
the two decisive witnesses are fixed and no voter-pair choice is needed.  With
three voters, the encoding must select which pair supplies the two witnesses,
making role-sensitive selector machinery a central part of the comparison.

\paragraph{Contributions.}
The paper makes three narrow contributions.
\begin{enumerate}
  \item It states a family-scale witness discipline that separates status
  agreement, clause-multiset equivalence, and repair-family comparison.
  \item It records a positive transfer along the two-voter alternative
  dimension: at $(2,5)$, clause-multiset equivalence persists under the identity
  map and raw repair non-canonicity persists.
  \item It records a voter-dimension boundary: at $(3,4)$, bundled and split
  packages are both UNSAT, but role-specific selector machinery breaks
  clause-multiset equivalence.  This case is classified
  \texttt{sat\_equiv\_only}, and no repair enumeration is authorized from that
  result.
\end{enumerate}

The result is intentionally modest.  We do not claim a theorem for arbitrary
$n$ and $m$, a new SAT algorithm, or a new encoding of full social acyclicity.
The evidence is restricted to the local clause families
\texttt{no\_cycle3} and \texttt{no\_cycle4}.  At five alternatives these
families do not exclude a length-five cycle.  The paper therefore reports a
controlled finite transfer study, not a universal social-choice theorem.

\paragraph{Roadmap.}
Section~\ref{sec:background} separates the semantic, encoding, and repair
layers.  Section~\ref{sec:witness} defines the comparison controls.
Sections~\ref{sec:option-d} and~\ref{sec:option-c} report the positive and
boundary tracks.  Section~\ref{sec:related} positions the result relative to
computational social choice, diagnosis, and certified SAT\@.
Section~\ref{sec:discussion} gives the combined interpretation and limitations.
Appendix~\ref{sec:repro} records the repository-facing reproduction contract.
